\vfill\pagebreak
\section{Functions that call themselves}
\label{sec:functions:recursion}

A function can call any function, including itself. Such a function is
\textit{recursive}. It solves a problem by solving a smaller version of the same
problem, and it relies on the smallest version having an answer that needs no
further call. A countdown from 3 is ``say 3, then count down from 2'', and a
countdown from 0 is simply ``liftoff''.

\lstinputlisting[style=coloredverbatim, caption={A recursive countdown, \textit{examples/functions/countdown.yr}}, label=lst:functions:countdown]{examples/functions/countdown.yr}
\vspace{-10pt}%
\begin{lstlisting}[style=bashVerb]
3...
2...
1...
Liftoff!
\end{lstlisting}

Every recursive function has the same two parts, and \token{countdown} shows
both:

\begin{itemize}
\item a \textit{base case}, the problem small enough to answer directly, without
  calling the function again: here \token{n == 0}, on line~4;
\item a \textit{recursive case}, which does a little of the work and hands the
  rest to a call on a smaller problem: here, print \token{n}, then count down from
  \token{n - 1}, on line~8.
\end{itemize}

Each call has its own parameter. When \token{countdown(3)} calls
\token{countdown(2)}, the new call gets a new \token{n}, holding 2, and the
\token{n} of the first call still holds 3. Nothing is shared between them.

\subsection{Waiting for the answer}

In \token{countdown}, the recursive call is the last thing the function does. In
most recursive functions, the call's result is needed to finish the work. The
\textit{factorial} of a number $n$, written $n!$, is the product
$1 \times 2 \times \dots \times n$, with $0! = 1$. It is also $n$ times the
factorial of $n - 1$, and that sentence is a recursive function.

\lstinputlisting[style=coloredverbatim, caption={The factorial, \textit{examples/functions/factorial.yr}}, label=lst:functions:factorial]{examples/functions/factorial.yr}
\vspace{-10pt}%
\begin{lstlisting}[style=bashVerb, escapechar=@+]
@+\textcolor{teal!80}{alice@dev:\mtilde}@+$ ./factorial
n = 20
20! = 2432902008176640000
\end{lstlisting}

To compute \token{n * factorial(n - 1)} on line~7, a call must wait for the
result of the call it makes, which waits for the next one, and so on down to the
base case. Figure~\ref{fig:functions:factorial_calls} shows the calls made by
\token{factorial(3)}. The chain goes down until \token{n} is 0, whose call
answers 1 at once. Then the results travel back up, and each waiting call
multiplies the result it receives by its own \token{n}.

\begin{figure}[h]
  \centering
  \begin{adjustbox}{max width=\linewidth}
    \begin{tikzpicture}[node distance=10mm,
        call/.style={cfstep, minimum width=52mm, align=left, text width=48mm}]

      \node[cfterm] (main) {main: println(factorial(3))};
      \node[call, below=of main] (f3) {n = 3:\quad 3 * factorial(2)};
      \node[call, below=of f3] (f2) {n = 2:\quad 2 * factorial(1)};
      \node[call, below=of f2] (f1) {n = 1:\quad 1 * factorial(0)};
      \node[call, below=of f1] (f0) {n = 0:\quad 1};

      \foreach \upper/\lower in {main/f3, f3/f2, f2/f1, f1/f0} {
        \draw[cfflow] ($(\lower.north west) + (8mm, 10mm)$) --
          node[cflab, left] {calls} ($(\lower.north west) + (8mm, 0)$);
      }
      \foreach \upper/\lower/\value in {main/f3/6, f3/f2/2, f2/f1/1, f1/f0/1} {
        \draw[cfflow, dashed] ($(\lower.north east) + (-8mm, 0)$) --
          node[cflab, right] {returns \value} ($(\lower.north east) + (-8mm, 10mm)$);
      }
    \end{tikzpicture}
  \end{adjustbox}
  \caption{\label{fig:functions:factorial_calls} The calls made by
    \token{factorial(3)}. Each call waits for the one below it, and each has its
    own \token{n}. The call with \token{n = 0} answers without calling, and the
    results travel back up, each call multiplying by its own \token{n}.}
\end{figure}


The factorial grows fast. \token{factorial(20)} is the largest that fits in an
\token{i64}; \token{factorial(21)} overflows, and the program prints a negative
number, for the reasons of Section~\ref{sec:types:integer_overflow}.

Two functions can also call each other. The next two decide whether a positive
number is even or odd, knowing only that 0 is even: a number is even when the
number before it is odd, and odd when the number before it is even. Each
function is the recursive case of the other.

\begin{lstlisting}[style=coloredverbatim, caption=Two functions calling each other, label=lst:functions:parity]
use std::io;

fn isEven(n: i32)-> bool {
  if n == 0 { true } else { isOdd(n - 1) }
}

fn isOdd(n: i32)-> bool {
  if n == 0 { false } else { isEven(n - 1) }
}

fn main() {
  println(isEven(10), " ", isOdd(10));
}
\end{lstlisting}
\vspace{-10pt}%
\begin{lstlisting}[style=bashVerb]
true false
\end{lstlisting}

This is a slow way to test parity, which \token{n \% 2 == 0} answers in one
operation. It shows that the base case can live in any of the functions of the
cycle, as long as every chain of calls reaches one.

\subsection{Infinite recursion}
\label{sec:functions:infinite_recursion}

A loop whose condition never becomes false runs forever. A recursion whose base
case is never reached does worse: it crashes. The next countdown has lost its
base case, and every call makes another.

\lstinputlisting[style=coloredverbatim, caption={A recursion with no base case, \textit{examples/functions/runaway.yr}}, label=lst:functions:runaway]{examples/functions/runaway.yr}
\vspace{-10pt}%
\begin{lstlisting}[style=bashVerb, escapechar=@+]
@+\textcolor{teal!80}{alice@dev:\mtilde}@+$ ./runaway
3...
2...
1...
0...
-1...
@+\textit{(about 130\,000 more lines)}@+
-130822...
-130823...
Segmentation fault (core dumped)
\end{lstlisting}

The compiler accepts this program, and it seems at first to count forever, into
the negative numbers. Then it stops, and the last line, printed by the terminal
rather than by the program, says that the operating system killed it.

The reason is the waiting that Figure~\ref{fig:functions:factorial_calls} shows.
A call that has not returned yet keeps its parameters and its place in the
caller somewhere in memory, in a region called the \textit{call stack}. The
stack has a fixed size, a few megabytes, and each call takes a little of it.
Here no call ever returns, so the calls pile up until the stack is full, and the
next call has nowhere to go. This is a \textit{stack overflow}. The number of
calls it takes depends on the machine; about 130\,000 was the count on the
machine that ran this listing. Chapter~\ref{chap:layout} explains what the stack
holds.

A base case that exists but is never reached has the same effect. With the base
case of Listing~\ref{lst:functions:countdown}, the call \token{countdown(-1)}
tests \token{-1 == 0}, then \token{-2 == 0}, and so on: \token{n} moves away from
0 and never meets it. Writing the test as \token{n <= 0} closes the gap, and a
recursive function is safest when its base case catches every value that it
cannot reduce.

Two questions check a recursive function against both mistakes:

\begin{itemize}
\item Is there a case that returns without calling the function?
\item Does every recursive call get a problem strictly closer to that case,
  whatever the arguments?
\end{itemize}

\subsection{Too deep}

A recursion can also overflow the stack while being correct. The sum of the
numbers from 1 to $n$ is $n$ plus the sum from 1 to $n - 1$, and a function
written that way gives the right answer, as long as it is not asked for too
much.

\lstinputlisting[style=coloredverbatim, caption={A correct recursion that goes too deep, \textit{examples/functions/deep\_sum.yr}}, label=lst:functions:deep_sum]{examples/functions/deep_sum.yr}
\vspace{-10pt}%
\begin{lstlisting}[style=bashVerb, escapechar=@+]
@+\textcolor{teal!80}{alice@dev:\mtilde}@+$ ./deep_sum
500500
Segmentation fault (core dumped)
\end{lstlisting}

The sum up to 1\,000 needs 1\,001 calls waiting at the same time, and fits. The
sum up to 100\,000 needs 100\,001 of them, and overflows the stack before the
base case is reached. The loop of Listing~\ref{lst:control:for_sum} computes the
same sum for any $n$, with one call and one variable.

A recursion is a good fit when each call does something the loop version would
have to organize by hand, as in the next subsection or in the exercises. When
the recursion only counts down, one step at a time, a loop does the same work
without piling up calls.

\subsection{The Ackermann function}
\label{sec:functions:ackermann}

The Ackermann function, named after the mathematician Wilhelm Ackermann, takes
two numbers $m$ and $n$ and is defined by three rules:

\begin{itemize}
\item when $m$ is 0, the answer is $n + 1$;
\item when $n$ is 0, the answer is the value for $m - 1$ and $1$;
\item otherwise, the answer is the value for $m - 1$ and the value for $m$ and
  $n - 1$.
\end{itemize}

Each rule becomes one branch of the function. The third is unusual: the second
argument of the recursive call is itself a recursive call.

\lstinputlisting[style=coloredverbatim, caption={The Ackermann function, \textit{examples/functions/ackermann.yr}}, label=lst:functions:ackermann]{examples/functions/ackermann.yr}
\vspace{-10pt}%
\begin{lstlisting}[style=bashVerb]
1 2 3 4 5 6
2 3 4 5 6 7
3 5 7 9 11 13
5 13 29 61 125 253
\end{lstlisting}

The function always terminates. Every call either lowers $m$, or keeps $m$ and
lowers $n$, so the calls always reach $m = 0$ in the end. Yet the table shows
how fast the values grow. The row $m = 1$ is $n + 2$, the row $m = 2$ is
$2n + 3$, and the row $m = 3$ is already $2^{n+3} - 3$, which doubles at each
step. The row $m = 4$ is worse still:

\begin{itemize}
\item \token{ackermann(4, 1)} is 65\,533. Computing it takes almost three billion
  calls, some ten seconds on the machine that ran these listings, with 65\,536
  calls waiting at the deepest point.
\item \token{ackermann(4, 2)} is $2^{65536} - 3$, a number with 19\,729 digits.
  It does not fit in any integer type, and no computer could make the number of
  calls this function needs to compute it.
\end{itemize}

The Ackermann function is famous for a reason beyond its growth. Mathematicians
proved that it cannot be computed with \token{for} loops alone, whose number of
iterations is known when they start: recursion, or a loop whose end is not known
in advance, is needed. It is also a reminder that a function can be correct,
terminate, and still never finish in practice. How many calls a function makes
matters as much as whether it stops.
